\chapter{2001年京蒙皖卷（春理）}
\section{单选题}
\begin{problem}
集合 \(M = \{1, 2, 3, 4, 5\}\) 的子集个数是（\(\quad\)）

\choices
    {\(32\)}
    {\(31\)}
    {\(16\)}
    {\(15\)}
\end{problem}

\begin{answer}
A
\end{answer}

\begin{solution}
集合 \(M\) 有 \(5\) 个元素. 对每个元素，都有“取”或“不取”两种选择，所以子集总数为 \(2^{5}=32\). 因此选 A.
\end{solution}

\begin{problem}
函数 \(f(x)=a^x\)（\(a>0\)，\(a\neq1\)）对于任意的实数 \(x\)，\(y\) 都有（\(\quad\)）

\choices
    {\(f(xy) = f(x)f(y)\)}
    {\(f(xy) = f(x) + f(y)\)}
    {\(f(x + y) = f(x)f(y)\)}
    {\(f(x + y) = f(x) + f(y)\)}
\end{problem}

\begin{answer}
C
\end{answer}

\begin{solution}
由指数运算律，对任意实数 \(x\)，\(y\) 都有
\[
f(x+y)=a^{x+y}=a^xa^y=f(x)f(y)
\]
其余选项均可排除：选 A 时取 \(x=2\)，\(y=3\)，则左边为 \(a^6\)，右边为 \(a^5\)，不相等；选 B 或 D 时取 \(x=y=0\)，则左边为 \(1\)，右边为 \(2\). 因此选 C.
\end{solution}

\begin{problem}
\(\displaystyle\lim_{n\to\infty}\frac{\mathrm C_{2n}^{n}}{\mathrm C_{2n+2}^{n+1}}=\)（\(\quad\)）

\choices
    {\(0\)}
    {\(2\)}
    {\( \frac{1}{2} \)}
    {\(\frac{1}{4}\)}
\end{problem}

\begin{answer}
D
\end{answer}

\begin{solution}
由组合数的阶乘表达式，
\[
\frac{\mathrm C_{2n}^{n}}{\mathrm C_{2n+2}^{n+1}}=\frac{(2n)!}{(n!)^2}\cdot\frac{((n+1)!)^2}{(2n+2)!}=\frac{n+1}{2(2n+1)}
\]
于是
\[
\lim_{n\to\infty}\frac{\mathrm C_{2n}^{n}}{\mathrm C_{2n+2}^{n+1}}=\lim_{n\to\infty}\frac{n+1}{2(2n+1)}=\frac14
\]
所以选 D.
\end{solution}

\begin{problem}
函数 \(y=-\sqrt{1-x}\)（\(x\leqslant1\)）的反函数是（\(\quad\)）

\choices
    {\(y=x^{2}-1\ \left(-1\leqslant x\leqslant0\right)\)}
    {\(y=x^{2}-1\ \left(0\leqslant x\leqslant1\right)\)}
    {\(y=1-x^{2}\ \left(x\leqslant0\right)\)}
    {\(y=1-x^{2}\ \left(0\leqslant x\leqslant1\right)\)}
\end{problem}

\begin{answer}
C
\end{answer}

\begin{solution}
原函数的定义域为 \(x\leqslant1\)，且 \(y=-\sqrt{1-x}\leqslant0\)，所以其值域为 \((-\infty,0]\). 又
\[
y=-\sqrt{1-x}\Longleftrightarrow y^2=1-x\Longleftrightarrow x=1-y^2
\]
交换自变量与因变量，反函数为 \(y=1-x^2\)，定义域为原函数的值域 \(x\leqslant0\). 因此选 C.
\end{solution}

\begin{problem}
极坐标系中，圆 \( \rho = 4 \cos \theta + 3 \sin \theta \) 的圆心的坐标是（\(\quad\)）

\choices
    {\(\left(\frac{5}{2}, \arcsin \frac{3}{5}\right)\)}
    {\(\left(5, \arcsin \frac{4}{5}\right)\)}
    {\(\left(5, \arcsin \frac{3}{5}\right)\)}
    {\(\left(\frac{5}{2}, \arcsin \frac{4}{5}\right)\)}
\end{problem}

\begin{answer}
A
\end{answer}

\begin{solution}
利用 \(x=\rho\cos\theta\)，\(y=\rho\sin\theta\)，将极坐标方程化为直角坐标方程，得
\[
x^2+y^2=4x+3y\Longleftrightarrow
(x-2)^2+\left(y-\frac32\right)^2=\frac{25}{4}
\]
所以圆心的直角坐标为 \(\left(2,\frac32\right)\)，极径为
\[
\rho=\sqrt{2^2+\left(\frac32\right)^2}=\frac52
\]
圆心在第一象限，故极角满足 \(\sin\theta=\frac{3/2}{5/2}=\frac35\)，圆心的极坐标为 \(\left(\frac52,\arcsin\frac35\right)\). 因此选 A.
\end{solution}

\begin{problem}
设动点 \(P\) 在直线 \(x = 1\) 上，\(O\) 为坐标原点. 以 \(OP\) 为直角边，点 \(O\) 为直角顶点作等腰 \(\mathrm{Rt}\triangle OPQ\)，则动点 \(Q\) 的轨迹是（\(\quad\)）

\choices
    {圆}
    {两条平行直线}
    {抛物线}
    {双曲线}
\end{problem}

\begin{answer}
B
\end{answer}

\begin{solution}
设 \(P=(1,t)\)，其中 \(t\in\R\)，则 \(\overrightarrow{OP}=(1,t)\). 以 \(OP\) 为直角边、\(O\) 为直角顶点作等腰直角三角形时，\(\overrightarrow{OQ}\) 是 \(\overrightarrow{OP}\) 旋转 \(90^\circ\) 所得的向量，因此
\[
Q=(-t,1)\quad\text{或}\quad Q=(t,-1)
\]
且 \(t\) 可取任意实数，所以 \(Q\) 的轨迹是两条直线 \(y=1\) 与 \(y=-1\)，选 B.
\end{solution}

\begin{problem}
已知 \( f(x^{6}) = \log_2 x \)，那么 \( f(8) \) 等于（\(\quad\)）

\choices
    {\( \frac{4}{3} \)}
    {\(8\)}
    {\(18\)}
    {\( \frac{1}{2} \)}
\end{problem}

\begin{answer}
D
\end{answer}

\begin{solution}
由对数的定义域知 \(x>0\). 令 \(u=x^6\)，则 \(u>0\)，且 \(x=u^{1/6}\)，因而
\[
f(u)=\log_2u^{1/6}=\frac16\log_2u
\]
取 \(u=8\)，得
\[
f(8)=\frac16\log_2 8=\frac16\cdot3=\frac12
\]
所以选 D.
\end{solution}

\begin{problem}
若 \(A\)，\(B\) 是锐角 \(\triangle ABC\) 的两个内角，则点 \(P(\cos B - \sin A, \sin B - \cos A)\) 在（\(\quad\)）

\choices
    {第一象限}
    {第二象限}
    {第三象限}
    {第四象限}
\end{problem}

\begin{answer}
B
\end{answer}

\begin{solution}
因为 \(\triangle ABC\) 是锐角三角形，所以 \(A+B=\pi-C>\frac\pi2\). 因而 \(A>\frac\pi2-B\)，在 \(\left(0,\frac\pi2\right)\) 内正弦函数递增，故
\[
\sin A>\sin\left(\frac\pi2-B\right)=\cos B
\]
所以 \(P\) 的横坐标为负. 由于 \(B>\frac\pi2-A\)，且正弦函数在 \(\left(0,\frac\pi2\right)\) 内递增，有
\[
\sin B>\sin\left(\frac\pi2-A\right)=\cos A
\]
所以 \(P\) 的纵坐标为正. 因此 \(P\) 在第二象限，选 B.
\end{solution}

\begin{problem}
如果圆锥的侧面展开图是半圆，那么这个圆锥的顶角（圆锥轴截面中两条母线的夹角）是（\(\quad\)）

\choices
    {\(30^{\circ}\)}
    {\(45^{\circ}\)}
    {\(60^{\circ}\)}
    {\(90^{\circ}\)}
\end{problem}

\begin{answer}
C
\end{answer}

\begin{solution}
设圆锥底面半径为 \(r\)，母线长为 \(l\). 侧面展开图是半径为 \(l\) 的半圆，其弧长等于圆锥底面周长，所以
\[
\pi l=2\pi r\Longrightarrow l=2r
\]
设圆锥轴截面顶角为 \(\varphi\). 在轴截面的一半中，斜边为 \(l\)，与顶角对应的对边为 \(r\)，故
\[
\sin\frac\varphi2=\frac rl=\frac12
\]
由于 \(0<\frac\varphi2<\frac\pi2\)，所以 \(\frac\varphi2=30^\circ\)，即 \(\varphi=60^\circ\). 因此选 C.
\end{solution}

\begin{problem}
若实数 \(a\)，\(b\) 满足 \(a + b = 2\)，则 \(3^a + 3^b\) 的最小值是（\(\quad\)）

\choices
    {\(18\)}
    {\(6\)}
    {\(2\sqrt{3}\)}
    {\(2\sqrt[4]{3}\)}
\end{problem}

\begin{answer}
B
\end{answer}

\begin{solution}
因为 \(3^a\)，\(3^b\) 均为正数，由基本不等式，
\[
3^a+3^b\geqslant2\sqrt{3^a3^b}=2\cdot3^{(a+b)/2}=6
\]
当且仅当 \(3^a=3^b\) 时取等号；由于 \(3>1\)，这等价于 \(a=b\)，再结合 \(a+b=2\) 得 \(a=b=1\). 所以最小值为 \(6\)，选 B.
\end{solution}

\begin{problem}
如图是正方体的平面展开图.

在这个正方体中，
\begin{circlelist}
\item \(BM\) 与 \(ED\) 平行；
\item \(CN\) 与 \(BE\) 是异面直线；
\item \(CN\) 与 \(BM\) 成 \(60^{\circ}\) 角；
\item \(DM\) 与 \(BN\) 垂直.
\end{circlelist}
以上四个命题中，正确命题的序号是（\(\quad\)）

\problemresetlayout
\bitmapfigure[max width=0.36\linewidth]{img/2001/jing_meng_wan_science_spring/q11_fig1.png}

\choices
    {\(\circled{1}\)\(\circled{2}\)\(\circled{3}\)}
    {\(\circled{2}\)\(\circled{4}\)}
    {\(\circled{3}\)\(\circled{4}\)}
    {\(\circled{2}\)\(\circled{3}\)\(\circled{4}\)}
\end{problem}

\begin{answer}
C
\end{answer}

\begin{solution}
取正方体棱长为 \(1\)，以图中中央正方形 \(ABCD\) 所在平面为 \(z=0\)，令
\(A=(0,0,0)\)，\(B=(1,0,0)\)，\(C=(1,1,0)\)，\(D=(0,1,0)\).
由展开图折成立方体后，可取 \(E=(0,0,1)\)，\(N=(0,1,1)\)，\(M=(1,1,1)\). 于是
\(\overrightarrow{BM}=(0,1,1)\)，\(\overrightarrow{ED}=(0,1,-1)\)
二者不平行，命题\(\circled{1}\)错误；又
\[
\overrightarrow{CN}=(-1,0,1)=\overrightarrow{BE}
\]
所以 \(CN\parallel BE\)，命题\(\circled{2}\)错误.

此外，
\(\overrightarrow{CN}\cdot\overrightarrow{BM}=1\)，\(|CN|=|BM|=\sqrt2\)
故 \(CN\) 与 \(BM\) 所成的角为 \(60^\circ\)，命题\(\circled{3}\)正确. 最后
\(\overrightarrow{DM}=(1,0,1)\)，\(\overrightarrow{BN}=(-1,1,1)\)
且 \(\overrightarrow{DM}\cdot\overrightarrow{BN}=0\)，所以命题\(\circled{4}\)正确. 因此正确命题的序号为 \(\circled{3}\)，\(\circled{4}\)，选 C.
\end{solution}

\begin{problem}
根据市场调查结果，预测某种家用商品从年初开始的 \(n\) 个月内累积的需求量 \(S_{n}\)（\(\text{万件}\)）近似地满足 \(S_{n} = \frac{n}{90}(21n - n^{2} - 5)\)（\(n = 1\)，\(2\)，\(\cdots\)，\(12\)）. 按此预测，在本年度内，需求量超过 \(1.5\text{ 万件}\) 的月份是（\(\quad\)）

\choices
    {\(5\text{ 月}\)，\(6\text{ 月}\)}
    {\(6\text{ 月}\)，\(7\text{ 月}\)}
    {\(7\text{ 月}\)，\(8\text{ 月}\)}
    {\(8\text{ 月}\)，\(9\text{ 月}\)}
\end{problem}

\begin{answer}
C
\end{answer}

\begin{solution}
记第 \(n\) 个月的实际需求量为 \(T_n\)，并令 \(S_0=0\)，则 \(T_n=S_n-S_{n-1}\). 由题给式子，
\[
T_n=\frac{-n^2+15n-9}{30}
\]
需求量超过 \(1.5\) 万件等价于
\[
\frac{-n^2+15n-9}{30}>\frac32\quad\Longleftrightarrow\quad n^2-15n+54<0
\]
也就是 \((n-6)(n-9)<0\). 因 \(n=1\)，\(2\)，\(\ldots\)，\(12\)，可得 \(n=7\)，\(8\). 所以需求量超过 \(1.5\) 万件的月份是 \(7\) 月、\(8\) 月，选 C.
\end{solution}

\section{填空题}
\begin{problem}
已知球内接正方体的表面积为 \(S\)，那么球体积等于\fillinblank{}.
\end{problem}

\begin{answer}
\(\dfrac{\pi S\sqrt{2S}}{24}\)
\end{answer}

\begin{solution}
设正方体棱长为 \(s\). 由表面积条件，\(6s^2=S\). 球的直径等于正方体体对角线 \(s\sqrt3\)，所以球半径 \(r=\frac{\sqrt3}{2}s\)，从而 \(r^2=\frac S8\). 因此球体积为
\[
V=\frac43\pi r^3=\frac43\pi\left(\frac S8\right)^{3/2}=\frac{\pi S\sqrt{2S}}{24}
\]
\end{solution}

\begin{problem}
椭圆 \( x^{2} + 4y^{2} = 4 \) 长轴上一个顶点为 \(A\)，以 \(A\) 为直角顶点作一个内接于椭圆的等腰直角三角形，该三角形的面积是\fillinblank{}.
\end{problem}

\begin{answer}
\(\dfrac{16}{25}\)
\end{answer}

\begin{solution}
椭圆的长轴端点为 \((\pm2,0)\). 不妨取 \(A=(2,0)\)，并设另一个顶点为 \(B=(2+u,v)\). 取 \(AC\) 为 \(AB\) 逆时针旋转 \(90^\circ\) 后得到的向量，则 \(C=(2-v,u)\)；顺时针情形关于 \(x\) 轴对称，面积相同.

由 \(B\)，\(C\) 在椭圆上，分别得到
\(u^2+4v^2+4u=0\)，\(4u^2+v^2-4v=0\).
用前式乘以 \(4\) 减去后式，得 \(15v^2+16u+4v=0\)，即 \(u=-\frac{v(15v+4)}{16}\). 代回前式并因式分解，得
\[
v(5v-4)(45v^2+60v+64)=0
\]
由于 \(45v^2+60v+64\) 的判别式小于 \(0\)，且 \(v=0\) 给出退化三角形，所以 \(v=\frac45\)，进而 \(u=-\frac45\). 因此直角边 \(AB\) 的平方为
\[
AB^2=u^2+v^2=\frac{32}{25}
\]
等腰直角三角形面积为 \(\frac12AB^2=\frac{16}{25}\).
\end{solution}

\begin{problem}
已知 \(\sin^2\alpha + \sin^2\beta + \sin^2\gamma = 1\)（\(\alpha\)，\(\beta\)，\(\gamma\) 均为锐角）. 那么 \(\cos\alpha\cos\beta\cos\gamma\) 的最大值等于\fillinblank{}.
\end{problem}

\begin{answer}
\(\dfrac{2\sqrt6}{9}\)
\end{answer}

\begin{solution}
令 \(x=\sin^2\alpha\)，\(y=\sin^2\beta\)，\(z=\sin^2\gamma\)，则 \(x+y+z=1\)，且 \(x\)，\(y\)，\(z\) 均为正数. 又
\[
\cos^2\alpha\cos^2\beta\cos^2\gamma=(1-x)(1-y)(1-z)=(y+z)(z+x)(x+y)
\]
三个正数 \(x+y\)，\(y+z\)，\(z+x\) 的和为 \(2\)，由基本不等式，
\[
(x+y)(y+z)(z+x)\leqslant\left(\frac23\right)^3
\]
当 \(x=y=z=\frac13\) 时取等号，且对应的角均为锐角. 因此
\[
\cos\alpha\cos\beta\cos\gamma\leqslant\left(\frac23\right)^{3/2}=\frac{2\sqrt6}{9}
\]
最大值为 \(\frac{2\sqrt6}{9}\).
\end{solution}

\begin{problem}
已知 \(m\)，\(n\) 是直线，\(\alpha\)，\(\beta\)，\(\gamma\) 是平面. 给出下列命题：
\begin{circlelist}
\item 若 \(\alpha \perp \beta\)，\(\alpha \cap \beta = m\)，\(n \perp m\)，则 \( n \perp \alpha \) 或 \( n \perp \beta \)；
\item 若 \(\alpha \parallel \beta\)，\(\alpha \cap \gamma = m\)，\(\beta \cap \gamma = n\)，则 \( m \parallel n \)；
\item 若 \( m \) 不垂直于 \( \alpha \)，则 \( m \) 不可能垂直于 \( \alpha \) 内的无数条直线；
\item 若 \(\alpha \cap \beta = m\)，\(n \parallel m\)，且 \(n \not\subset \alpha\)，\(n \not\subset \beta\)，则 \( n \parallel \alpha \) 且 \( n \parallel \beta \).
\end{circlelist}
其中正确的命题的序号是\fillinblank{}.
\end{problem}

\begin{answer}
\(\circled{2}\)，\(\circled{4}\)
\end{answer}

\begin{solution}
命题\(\circled{1}\)不成立. 例如取 \(\alpha:z=0\)，\(\beta:y=0\)，则 \(m\) 为 \(x\) 轴. 取过原点且方向向量为 \((0,1,1)\) 的直线 \(n\)，则 \(n\perp m\)，但 \(n\) 的方向既不平行于 \(\alpha\) 的法向量，也不平行于 \(\beta\) 的法向量，所以 \(n\) 不垂直于 \(\alpha\) 或 \(\beta\).

命题\(\circled{2}\)成立. 两个平行平面被同一平面 \(\gamma\) 所截时，截交线的方向相同，所以 \(m\parallel n\).

命题\(\circled{3}\)不成立. 若 \(m\) 在平面 \(\alpha\) 内且不垂直于 \(\alpha\)，则在 \(\alpha\) 内过 \(m\) 上各点分别作 \(m\) 的垂线，可得到无数条与 \(m\) 垂直的直线.

命题\(\circled{4}\)成立. 直线 \(n\) 平行于平面 \(\alpha\cap\beta\) 的直线 \(m\)，且 \(n\) 不在 \(\alpha\)，所以 \(n\parallel\alpha\)；同样，\(n\parallel m\)，\(m\subset\beta\)，而 \(n\not\subset\beta\)，故 \(n\parallel\beta\). 因此正确命题为 \(\circled{2}\)，\(\circled{4}\).
\end{solution}

\section{解答题}
\begin{problem}
设函数 \(f(x)=\frac{x+a}{x+b}\)（\(a>b>0\)），求 \(f(x)\) 的单调区间，并证明 \(f(x)\) 在其单调区间上的单调性.
\end{problem}

\begin{solution}
函数定义域为 \(x\neq-b\). 任取同一单调区间内的 \(x_1<x_2\)，则
\[
f(x_1)-f(x_2)=\frac{(x_1+a)(x_2+b)-(x_2+a)(x_1+b)}{(x_1+b)(x_2+b)}=\frac{(b-a)(x_1-x_2)}{(x_1+b)(x_2+b)}
\]
因 \(a>b\)，有 \(b-a<0\). 当 \(x_1\)，\(x_2\in(-\infty,-b)\) 或 \(x_1\)，\(x_2\in(-b,+\infty)\) 时，\((x_1+b)(x_2+b)>0\)，而 \(x_1-x_2<0\)，所以 \(f(x_1)-f(x_2)>0\)，即 \(f(x_1)>f(x_2)\). 因此 \(f(x)\) 在 \((-\infty,-b)\) 和 \((-b,+\infty)\) 上均单调递减.
\end{solution}

\begin{problem}
已知 \(z^{7}=1\)（\(z\in\mathbf{C}\)，\(z\neq1\)）.
\begin{enumerate}
\item 证明 \(1 + z + z^2 + z^3 + z^4 + z^5 + z^6 = 0\)；
\item 设 \(z\) 的辐角为 \(\alpha\)，求 \(\cos \alpha + \cos 2\alpha + \cos 4\alpha\) 的值.
\end{enumerate}
\end{problem}

\begin{solution}
由 \(z^7=1\) 且 \(z\neq1\)，有
\[
0=z^7-1=(z-1)(1+z+z^2+z^3+z^4+z^5+z^6)
\]
因 \(z-1\neq0\)，所以 \(1+z+z^2+z^3+z^4+z^5+z^6=0\).

又因 \(z\) 是单位复数，且 \(\arg z=\alpha\)，所以
\[
\cos\alpha+\cos2\alpha+\cos4\alpha=\mathop{\rm Re}(z+z^2+z^4)
\]
令 \(T=z+z^2+z^4\). 由 \(z^7=1\) 得 \(\overline z=z^6\)，于是
\[
\overline T=z^6+z^5+z^3
\]
利用前一问结论，\(T+\overline T=z+z^2+\cdots+z^6=-1\). 因而 \(2\mathop{\rm Re}T=-1\)，所求值为 \(-\frac12\).
\end{solution}

\begin{problem}
已知 \(VC\) 是 \(\triangle ABC\) 所在平面的一条斜线，点 \(N\) 是 \(V\) 在平面 \(ABC\) 上的射影，且在 \(\triangle ABC\) 的高 \(CD\) 上. \(AB = a\)，\(VC\) 与 \(AB\) 之间的距离为 \(h\)，点 \(M \in VC\).
\begin{enumerate}
\item 证明 \(\angle MDC\) 是二面角 \(M - AB - C\) 的平面角；
\item 当 \(\angle MDC = \angle CVN\) 时，证明 \(VC \perp\) 平面 \(AMB\)；
\item 若 \(\angle MDC = \angle CVN = \theta\)（\(0 < \theta < \frac{\pi}{2}\)），求四面体 \(MABC\) 的体积.
\end{enumerate}
\bitmapfigure[max width=0.34\linewidth]{img/2001/jing_meng_wan_science_spring/q19_fig1.png}
\end{problem}

\begin{solution}
建立直角坐标系，使 \(D=(0,0,0)\)，\(AB\) 为 \(x\) 轴，平面 \(ABC\) 为 \(z=0\). 记 \(CD=c\)，\(CN=r\)，\(VN=v\)，并取正方向使 \(C=(0,c,0)\)，\(N=(0,c-r,0)\)，\(V=(0,c-r,v)\)，其中 \(c\)，\(r\)，\(v>0\). 设 \(M\) 将 \(CV\) 按参数 \(t\) 分点，则
\[
M=(0,c-tr,tv)
\]

因为 \(D\)，\(A\)，\(B\) 共线，\(\overrightarrow{DC}=(0,c,0)\) 在平面 \(ABC\) 内且垂直于 \(AB\)，而 \(\overrightarrow{DM}=(0,c-tr,tv)\) 在平面 \(MAB\) 内且也垂直于 \(AB\). 所以 \(\angle MDC\) 是二面角 \(M-AB-C\) 的平面角.

在直角三角形 \(CVN\) 中，\(\tan\angle CVN=\frac rv\). 又
\[
\tan\angle MDC=\frac{tv}{c-tr}
\]
当两角相等时，\(\frac{tv}{c-tr}=\frac rv\)，即 \(tv^2=r(c-tr)\). 因而
\[
\overrightarrow{DM}\cdot\overrightarrow{CV}=-(c-tr)r+tv^2=0
\]
由坐标可得 \(\overrightarrow{CV}=(0,-r,v)\)，而 \(\overrightarrow{AB}\) 的方向向量为 \((1,0,0)\)，故 \(VC\perp AB\). 上式又给出 \(VC\perp DM\)，且 \(AB\) 与 \(DM\) 是平面 \(AMB\) 内相交的两条直线，所以 \(VC\perp\) 平面 \(AMB\).

令 \(L=VC=\sqrt{r^2+v^2}\). 由 \(\angle CVN=\theta\)，有 \(r=L\sin\theta\)，\(v=L\cos\theta\). 在垂直于 \(AB\) 的截面内，\(AB\) 变为过 \(D\) 的点，直线 \(VC\) 的方向向量为 \((-r,v)\)，由点到直线的距离公式，直线 \(VC\) 到 \(AB\) 的距离为
\[
h=\frac{cv}{L}=c\cos\theta
\]
故 \(c=\frac{h}{\cos\theta}\). 由 \(tv^2=r(c-tr)\) 得 \(tL^2=rc\)，从而点 \(M\) 到平面 \(ABC\) 的距离为
\[
tv=\frac{rcv}{L^2}=c\sin\theta\cos\theta
\]
于是
\[
V_{MABC}=\frac13\cdot\frac12ac\cdot c\sin\theta\cos\theta=\frac{ac^2\sin\theta\cos\theta}{6}=\frac{ah^2\tan\theta}{6}
\]
\end{solution}

\begin{problem}
在 \(1\) 与 \(2\) 之间插入 \(n\) 个正数 \(a_{1}\)，\(a_{2}\)，\(a_{3}\)，\(\cdots\)，\(a_{n}\)，使这 \(n+2\) 个数变成等比数列；又在 \(1\) 与 \(2\) 之间插入 \(n\) 个正数 \(b_{1}\)，\(b_{2}\)，\(b_{3}\)，\(\cdots\)，\(b_{n}\)，使这 \(n+2\) 个数成等差数列. 记 \(A_{n} = a_{1}a_{2}a_{3} \cdots a_{n}\)，\(B_{n} = b_{1} + b_{2} + b_{3} + \cdots + b_{n}\).
\begin{enumerate}
\item 求数列 \(\{A_{n}\}\) 和 \(\{B_{n}\}\) 的通项；
\item 当 \(n \geqslant 7\) 时，比较 \(A_{n}\) 与 \(B_{n}\) 的大小，并证明你的结论.
\end{enumerate}
\end{problem}

\begin{solution}
等比数列共有 \(n+2\) 项，首项为 \(1\)，末项为 \(2\)，设公比为 \(q\). 则 \(q^{n+1}=2\)，所以
\(a_k=q^k=2^{\frac{k}{n+1}}\)，\((k=1,2,\ldots,n)\).
因此
\[
A_n=q^{1+2+\cdots+n}=2^{\frac{n(n+1)}{2(n+1)}}=2^{n/2}
\]

等差数列的公差为 \(\frac1{n+1}\)，所以
\(b_k=1+\frac{k}{n+1}\)，\((k=1,2,\ldots,n)\)
从而
\[
B_n=\sum_{k=1}^n\left(1+\frac{k}{n+1}\right)=n+\frac n2=\frac{3n}{2}
\]

当 \(n=7\) 时，\(2^n=128>\frac94n^2=\frac94\cdot49=\frac{441}{4}\). 对 \(n\geqslant3\)，有
\[
\frac{2^{n+1}/(n+1)^2}{2^n/n^2}=\frac{2n^2}{(n+1)^2}>1
\]
所以 \(\frac{2^n}{n^2}\) 在 \(n\geqslant3\) 时递增. 因而对所有 \(n\geqslant7\)，都有 \(2^n>\frac94n^2\)，即 \(2^{n/2}>\frac{3n}{2}\). 所以 \(A_n>B_n\).
\end{solution}

\begin{problem}
某摩托车生产企业，上年度生产摩托车的投入成本为 \(1\text{ 万元/辆}\)，出厂价为 \(1.2\text{ 万元/辆}\)，年销售量为 \(1000\text{ 辆}\). 本年度为适应市场需求，计划提高产品档次，适度增加投入成本. 若每辆车投入成本增加的比例为 \(x\)（\(0 < x < 1\)），则出厂价相应提高的比例为 \(0.75x\). 同时预计年销售量增加的比例为 \(0.6x\). 已知年利润 \(= (\text{出厂价} - \text{投入成本}) \times \text{年销售量}\).
\begin{enumerate}
\item 写出本年度预计的年利润 \(y\) 与投入成本增加的比例 \(x\) 的关系式；
\item 为使本年度的年利润比上年有所增加，问投入成本增加的比例 \(x\) 应在什么范围内？
\end{enumerate}
\end{problem}

\begin{solution}
本年度每辆车的投入成本为 \(1+x\) 万元，出厂价为
\[
1.2(1+0.75x)=1.2+0.9x
\]
万元，年销售量为 \(1000(1+0.6x)\) 辆. 因此预计年利润为
\[
y=(1.2+0.9x-1-x)\cdot1000(1+0.6x)=200+20x-60x^2
\]
其中 \(0<x<1\).

上年度利润为 \((1.2-1)\times1000=200\) 万元. 要使本年度利润增加，需
\[
200+20x-60x^2>200
\]
即 \(20x(1-3x)>0\). 结合 \(0<x<1\)，得到 \(0<x<\frac13\).
\end{solution}

\begin{problem}
已知抛物线 \(y^{2} = 2px\)（\(p > 0\)）. 过动点 \(M(a,0)\) 且斜率为 1 的直线 \(l\) 与该抛物线交于不同的两点 \(A\)，\(B\)，\(|AB| \leqslant 2p\).
\begin{enumerate}
\item 求 \(a\) 的取值范围；
\item 若线段 \(AB\) 的垂直平分线交 \(x\) 轴于点 \(N\)，求 \(\mathrm{Rt}\triangle NAB\) 面积的最大值.
\end{enumerate}
\end{problem}

\begin{solution}
直线 \(l\) 的方程为 \(y=x-a\)，即 \(x=y+a\). 设交点 \(A\)，\(B\) 的纵坐标分别为 \(y_1\)，\(y_2\)，代入抛物线方程得
\(y^2=2p(y+a)\)，\(y^2-2py-2pa=0\).
两交点不同等价于判别式 \(4p(p+2a)>0\)，所以 \(a>-\frac p2\). 又由韦达定理，
\[
|y_1-y_2|=2\sqrt{p(p+2a)}
\]
由于直线 \(l\) 的斜率为 \(1\)，有
\[
|AB|=\sqrt{(x_1-x_2)^2+(y_1-y_2)^2}=\sqrt2|y_1-y_2|=2\sqrt{2p(p+2a)}
\]
由 \(|AB|\leqslant2p\)，得 \(a\leqslant-\frac p4\). 因此
\[
-\frac p2<a\leqslant-\frac p4
\]

设 \(Q\) 为线段 \(AB\) 的中点. 由 \(y_1+y_2=2p\)，得
\[
Q=(a+p,p)
\]
线段 \(AB\) 的斜率为 \(1\)，所以其垂直平分线斜率为 \(-1\)，方程为 \(x+y=a+2p\). 它与 \(x\) 轴交于 \(N=(a+2p,0)\)，而 \(M=(a,0)\). 因此
\(\overrightarrow{QM}=(-p,-p)\)，\(\overrightarrow{QN}=(p,-p)\)
两向量垂直且长度均为 \(p\sqrt2\)，所以 \(NQ=p\sqrt2\)，且 \(NQ\perp AB\). 于是
\[
S_{\triangle NAB}=\frac12|AB|\cdot NQ\leqslant\frac12\cdot2p\cdot p\sqrt2=\sqrt2p^2
\]
当 \(a=-\frac p4\) 时 \(|AB|=2p\)，等号成立. 所以面积的最大值为 \(\sqrt2p^2\).
\end{solution}
